% GENERATED FILE. Edit canonical lessons, then run npm run generate:books.
% canonical-source-hash: fnv1a64:bdc2e713d41ce48b
% canonical-lessons: HI-C89-kar-raha-hun, HI-C89-kar-rahi-hun, HI-C89-kar-raha-hai, HI-C89-kar-rahe-hain, HI-C89-karta-vs-kar-raha, HI-C89-repaso-jaari, HI-C89-sintesis-jaari

\chapter{What Is Going On Right Now}
\label{ch:hi-jaari}
\hlchaptermodality{97}

\begin{chapteropening}
\textbf{By the end of this chapter:} \emph{I can say what I am doing at this moment, ask somebody else what they are doing, and keep that apart from what I do in general.}

Ask what somebody is doing at the aap level and answer for yourself, then say the same thing as a habit rather than as a moment.
\end{chapteropening}

\section[kar rahā hū̃]{\hi{कर} \hi{रहा} \hi{हूँ} — one word slid into the middle}

\label{lesson:HI-C89-kar-raha-hun}



\begin{quote}
Say \emph{I do} — the habit. Now try to say what you are doing \textbf{at this moment}, and notice the habit form does not quite mean it.
\end{quote}

\subsection*{You'll want to know: \hi{कर} \hi{रहा} \hi{हूँ}}
\textbf{\hi{कर} \hi{रहा} \hi{हूँ}} (\emph{kar rahā hū̃}) — \textquotedblleft{}I am doing,\textquotedblright{} right now, said by a man.

Three words: the bare stem, then \textbf{\hi{रहा}}, then the copula you have had since your third chapter.

\begin{cousinweb}
\noindent
\begin{tabularx}{\linewidth}{@{}>{\raggedright\arraybackslash}X>{\raggedright\arraybackslash}X@{}}
\toprule
\textbf{piece} & \textbf{what it carries} \\
\midrule
\hi{कर} & the action, bare \\
\hi{रहा} & that it is in progress \\
\hi{हूँ} & who, and that it is now \\
\bottomrule
\end{tabularx}

\textbf{\hi{रहा}} is the past participle of \textbf{\hi{रहना}}, \textquotedblleft{}to stay, to remain.\textquotedblright{} So the sentence says, almost literally, \emph{doing — having stayed — I am}. The action is pictured as one you are \textbf{still inside}.

English does the same thing from a different angle: \emph{I am doing} uses a participle too, and \emph{I remain at work} is the same metaphor spelled out.
\end{cousinweb}

\begin{grammarlens}[title={Grammar Lens: the stem is bare, and that is the point}]
Look at what happens to the verb in each tense:

\noindent
\begin{tabularx}{\linewidth}{@{}>{\raggedright\arraybackslash}X>{\raggedright\arraybackslash}X@{}}
\toprule
\textbf{what you say} & \textbf{the verb} \\
\midrule
I do (habit) & \hi{करता} \\
I am doing (now) & \hi{कर} \\
\bottomrule
\end{tabularx}

The habit puts an ending on the stem. The continuous \textbf{takes the ending off} and puts a separate word after it instead. Nothing else in the sentence moves.

That means the pattern works on \textbf{any} verb you know. Take the stem, add \textbf{\hi{रहा}}, keep the copula: \emph{\hi{बोल} \hi{रहा} \hi{हूँ}}, \emph{\hi{जा} \hi{रहा} \hi{हूँ}}, \emph{\hi{खा} \hi{रहा} \hi{हूँ}}.
\end{grammarlens}

\subsection*{Your turn}
Say these aloud:

\begin{itemize}
\raggedright
  \item \textquotedblleft{}\hi{कर} \hi{रहा} \hi{हूँ}\textquotedblright{} — I am doing
  \item \textquotedblleft{}\hi{बोल} \hi{रहा} \hi{हूँ}\textquotedblright{}
  \item the three pieces, and what each one carries
  \item what \hi{रहना} means on its own
\end{itemize}

\begin{tcolorbox}[breakable,colback=teal!4,colframe=teal!35!black,title={Before you move on}]
Which word carries \emph{in progress}? (\textbf{\hi{रहा}}.) What shape is the verb stem in? (Bare — the habit ending comes off.) What does \textbf{\hi{रहना}} mean by itself? (To stay, to remain.)
\end{tcolorbox}
\section[kar rahī hū̃]{\hi{कर} \hi{रही} \hi{हूँ} — which of the three words listens}

\label{lesson:HI-C89-kar-rahi-hun}



\begin{quote}
\textbf{\hi{रहा}} ends in \textbf{-\hi{आ}}. Before reading on, guess what it does when a woman is speaking, and guess which of the other two words joins in.
\end{quote}

\subsection*{You'll want to know: \hi{कर} \hi{रही} \hi{हूँ}}
\textbf{\hi{कर} \hi{रही} \hi{हूँ}} (\emph{kar rahī hū̃}) — \textquotedblleft{}I am doing,\textquotedblright{} said by a woman.

One letter has changed, and it is in the middle word.

\begin{cousinweb}
\noindent
\begin{tabularx}{\linewidth}{@{}>{\raggedright\arraybackslash}X>{\raggedright\arraybackslash}X@{}}
\toprule
\textbf{speaker} & \textbf{the three words} \\
\midrule
a man & \hi{कर} \hi{रहा} \hi{हूँ} \\
a woman & \hi{कर} \hi{रही} \hi{हूँ} \\
\bottomrule
\end{tabularx}

The stem stays bare either way. The copula \textbf{\hi{हूँ}} is about \emph{who is speaking} rather than about gender, so it does not move. Only \textbf{\hi{रहा}} bends, and it bends the way every \textbf{-\hi{आ}} word you have met bends.

That is the tidiest arrangement you have seen yet: \textbf{one word per job, and only the one that describes the action listens to gender.}
\end{cousinweb}

\begin{grammarlens}[title={Grammar Lens: the same ending, a fourth time}]
\noindent
\begin{tabularx}{\linewidth}{@{}>{\raggedright\arraybackslash}X>{\raggedright\arraybackslash}X@{}}
\toprule
\textbf{masculine} & \textbf{feminine} \\
\midrule
\hi{काला} & \hi{काली} \\
\hi{बोलता} & \hi{बोलती} \\
\hi{था} & \hi{थी} \\
\hi{रहा} & \hi{रही} \\
\bottomrule
\end{tabularx}

Four places, one alternation. Every time a Hindi word ends in \textbf{-\hi{आ}} and describes somebody, \textbf{-\hi{ई}} is what it becomes for a woman — and by now you should expect it before you are told.
\end{grammarlens}

\subsection*{Your turn}
Say these aloud:

\begin{itemize}
\raggedright
  \item \textquotedblleft{}\hi{कर} \hi{रही} \hi{हूँ}\textquotedblright{} — I am doing
  \item both forms one after the other, and hear the one letter
  \item \textquotedblleft{}\hi{बोल} \hi{रही} \hi{हूँ}\textquotedblright{}
  \item which of the three words does not care about gender, and why
\end{itemize}

\begin{tcolorbox}[breakable,colback=teal!4,colframe=teal!35!black,title={Before you move on}]
Which of the three words bends? (The middle one, \textbf{\hi{रहा} / \hi{रही}}.) Why does the copula stay put? (It is about who, not about gender.) What is the fourth pair on the list? (\textbf{\hi{रहा} / \hi{रही}}.)
\end{tcolorbox}
\section[kar rahā hai]{\hi{कर} \hi{रहा} \hi{है} — two words doing two different jobs}

\label{lesson:HI-C89-kar-raha-hai}



\begin{quote}
Say \emph{I am doing}. Now point at somebody else and change only what has to change.
\end{quote}

\subsection*{You'll want to know: \hi{कर} \hi{रहा} \hi{है}}
\textbf{\hi{वह} \hi{कर} \hi{रहा} \hi{है।}} — He is doing (it). \textbf{\hi{वह} \hi{कर} \hi{रही} \hi{है।}} — She is doing (it).

\begin{cousinweb}
Only the \textbf{last} word moved when the subject changed from \emph{I} to \emph{he}:

\noindent
\begin{tabularx}{\linewidth}{@{}>{\raggedright\arraybackslash}X>{\raggedright\arraybackslash}X>{\raggedright\arraybackslash}X@{}}
\toprule
\textbf{subject} & \textbf{middle} & \textbf{last} \\
\midrule
\hi{मैं} & \hi{रहा} / \hi{रही} & \hi{हूँ} \\
\hi{वह} & \hi{रहा} / \hi{रही} & \hi{है} \\
\bottomrule
\end{tabularx}

Two independent dials. The middle word answers \emph{what is the subject like} and the last word answers \emph{which person is it}. Turning one does nothing to the other, which is why you can build any combination without a table to memorise.
\end{cousinweb}

\begin{grammarlens}[title={Grammar Lens: it works on things as well as people}]
\noindent
\begin{tabularx}{\linewidth}{@{}>{\raggedright\arraybackslash}X>{\raggedright\arraybackslash}X@{}}
\toprule
\textbf{sentence} & \textbf{meaning} \\
\midrule
\hi{वह} \hi{जा} \hi{रहा} \hi{है।} & He is going. \\
\hi{वह} \hi{आ} \hi{रही} \hi{है।} & She is coming. \\
\bottomrule
\end{tabularx}

A noun that is not a person still has a gender, and the middle word follows the noun rather than any person. That is how the weather sentence you learned long ago was put together — the rain is feminine, so the middle word ended in \textbf{-\hi{ई}} — and the review lesson takes it apart.
\end{grammarlens}

\subsection*{Your turn}
Say these aloud:

\begin{itemize}
\raggedright
  \item \textquotedblleft{}\hi{वह} \hi{कर} \hi{रहा} \hi{है}\textquotedblright{}
  \item the same about a woman
  \item which word changed when you moved from \emph{I} to \emph{he}
  \item \textquotedblleft{}\hi{वह} \hi{जा} \hi{रहा} \hi{है}\textquotedblright{}
\end{itemize}

\begin{tcolorbox}[breakable,colback=teal!4,colframe=teal!35!black,title={Before you move on}]
Which word answers \emph{which person}? (The last one.) Which answers \emph{what is the subject like}? (The middle one.) Do the two ever interfere? (No — they are separate dials.)
\end{tcolorbox}
\section[kar rahe haiṁ]{\hi{कर} \hi{रहे} \hi{हैं} — both dials turned to plural}

\label{lesson:HI-C89-kar-rahe-hain}



\begin{quote}
You have two dials: the middle word and the last. Say the pronoun for \emph{we}, and guess what a group does to each dial.
\end{quote}

\subsection*{You'll want to know: \hi{कर} \hi{रहे} \hi{हैं}}
\textbf{\hi{हम} \hi{कर} \hi{रहे} \hi{हैं।}} — We are doing (it). \textbf{\hi{आप} \hi{क्या} \hi{कर} \hi{रहे} \hi{हैं}?} — What are you doing?

\begin{cousinweb}
This time \textbf{both} words move:

\noindent
\begin{tabularx}{\linewidth}{@{}>{\raggedright\arraybackslash}X>{\raggedright\arraybackslash}X>{\raggedright\arraybackslash}X@{}}
\toprule
\textbf{subject} & \textbf{middle} & \textbf{last} \\
\midrule
\hi{मैं} & \hi{रहा} & \hi{हूँ} \\
\hi{वह} & \hi{रहा} & \hi{है} \\
\hi{हम}, \hi{वे}, \hi{आप} & \hi{रहे} & \hi{हैं} \\
\bottomrule
\end{tabularx}

\textbf{\hi{रहे}} is the plural of the same \textbf{-\hi{आ} / -\hi{ई} / -\hi{ए}} set. \textbf{\hi{हैं}} is the plural copula. Nothing here is new — it is the two dials you already have, both turned one notch.

At the \textbf{\hi{तुम}} level the last word is \textbf{\hi{हो}}: \emph{\hi{तुम} \hi{क्या} \hi{कर} \hi{रहे} \hi{हो}?}
\end{cousinweb}

\begin{grammarlens}[title={Grammar Lens: the question a listening paper asks}]
\noindent
\begin{tabularx}{\linewidth}{@{}>{\raggedright\arraybackslash}X>{\raggedright\arraybackslash}X@{}}
\toprule
\textbf{level} & \textbf{the question} \\
\midrule
\hi{आप} & \hi{आप} \hi{क्या} \hi{कर} \hi{रहे} \hi{हैं}? \\
\hi{तुम} & \hi{तुम} \hi{क्या} \hi{कर} \hi{रहे} \hi{हो}? \\
\bottomrule
\end{tabularx}

\emph{What are you doing?} is one of the commonest openings in spoken Hindi, and the whole difference between the two rows is the last word. That is the level decision again, settling one more system.
\end{grammarlens}

\subsection*{Your turn}
Say these aloud:

\begin{itemize}
\raggedright
  \item \textquotedblleft{}\hi{हम} \hi{कर} \hi{रहे} \hi{हैं}\textquotedblright{}
  \item \textquotedblleft{}\hi{आप} \hi{क्या} \hi{कर} \hi{रहे} \hi{हैं}?\textquotedblright{}
  \item the same question at the \emph{tum} level
  \item which two words moved, and which stayed bare
\end{itemize}

\begin{tcolorbox}[breakable,colback=teal!4,colframe=teal!35!black,title={Before you move on}]
What is the plural of \textbf{\hi{रहा}}? (\textbf{\hi{रहे}}.) Which word separates the \emph{āp} question from the \emph{tum} one? (The last — \textbf{\hi{हैं}} against \textbf{\hi{हो}}.) What shape is the stem in, still? (Bare.)
\end{tcolorbox}
\section[do vartamān]{do vartamān — the habit and the moment}

\label{lesson:HI-C89-karta-vs-kar-raha}



\begin{quote}
Say both forms of \emph{to do} that you now have — the habit and the one in progress — and try to say in English what separates them.
\end{quote}

\subsection*{You'll want to know: the pair}
\noindent
\begin{tabularx}{\linewidth}{@{}>{\raggedright\arraybackslash}X>{\raggedright\arraybackslash}X@{}}
\toprule
\textbf{Hindi} & \textbf{English} \\
\midrule
\hi{मैं} \hi{हिंदी} \hi{बोलता} \hi{हूँ।} & I speak Hindi. \\
\hi{मैं} \hi{हिंदी} \hi{बोल} \hi{रहा} \hi{हूँ।} & I am speaking Hindi. \\
\bottomrule
\end{tabularx}

The first is a \textbf{fact about you}. The second is a \textbf{report about this minute}.

\begin{cousinweb}
This distinction is worth pausing on, because English speakers have it and most European learners of Hindi do not. French, German and Spanish all translate both of those lines with a single present tense; Hindi and English each keep two.

So the good news is that \textbf{your instinct is already right}. Where English says \emph{I am doing}, Hindi says \emph{\hi{कर} \hi{रहा} \hi{हूँ}}. Where English says \emph{I do}, Hindi says \emph{\hi{करता} \hi{हूँ}}. Translate the shape, not the words.

One caution. Hindi leans on the habit form in a place English does not: for something happening in the near future that is already arranged, Hindi is happy with the habitual, where English reaches for \emph{I am}.
\end{cousinweb}

\begin{grammarlens}[title={Grammar Lens: which word in the sentence tells you}]
\noindent
\begin{tabularx}{\linewidth}{@{}>{\raggedright\arraybackslash}X>{\raggedright\arraybackslash}X@{}}
\toprule
\textbf{form} & \textbf{the shape of the verb} \\
\midrule
habit & stem \textbf{+ -\hi{ता}}, then the copula \\
in progress & bare stem, then \textbf{\hi{रहा}}, then the copula \\
\bottomrule
\end{tabularx}

The copula is in both, so it tells you nothing. What decides it is whether the stem carries \textbf{-\hi{ता}} or stands bare with \textbf{\hi{रहा}} behind it.

\textbf{\hi{अभी}} — the word for \emph{right now} — sits comfortably with the second and strangely with the first, which is the quickest way to hear the difference.
\end{grammarlens}

\subsection*{Your turn}
Say these aloud:

\begin{itemize}
\raggedright
  \item \textquotedblleft{}\hi{मैं} \hi{हिंदी} \hi{बोलता} \hi{हूँ}\textquotedblright{} — the fact
  \item \textquotedblleft{}\hi{मैं} \hi{हिंदी} \hi{बोल} \hi{रहा} \hi{हूँ}\textquotedblright{} — the minute
  \item which one you would use to answer \emph{what are you doing?}
  \item the one that takes \textbf{\hi{अभी}} without sounding odd
\end{itemize}

\begin{tcolorbox}[breakable,colback=teal!4,colframe=teal!35!black,title={Before you move on}]
Which form states a fact about you? (The habit, with \textbf{-\hi{ता}}.) Which reports this minute? (The one with \textbf{\hi{रहा}}.) Does the copula help you tell them apart? (No — it is in both.)
\end{tcolorbox}
\section[jārī kām kī tālikā]{jārī kām kī tālikā — two dials, and two sentences you already knew}

\label{lesson:HI-C89-repaso-jaari}



\begin{quote}
Name the three pieces of a continuous sentence, and say which two of them move, before you look at the table.
\end{quote}

\begin{grammarlens}[title={Grammar Lens: the grid}]
\noindent
\begin{tabularx}{\linewidth}{@{}>{\raggedright\arraybackslash}X>{\raggedright\arraybackslash}X@{}}
\toprule
\textbf{subject} & \textbf{middle and last} \\
\midrule
\hi{मैं} (m) & \hi{रहा} \hi{हूँ} \\
\hi{मैं} (f) & \hi{रही} \hi{हूँ} \\
\hi{वह} (m) & \hi{रहा} \hi{है} \\
\hi{वह} (f) & \hi{रही} \hi{है} \\
\hi{हम}, \hi{वे}, \hi{आप} & \hi{रहे} \hi{हैं} \\
\hi{तुम} & \hi{रहे} \hi{हो} \\
\bottomrule
\end{tabularx}

The stem in front of all of them is \textbf{bare}, every time. Six rows and not one of them touches the verb.
\end{grammarlens}

\begin{grammarlens}[title={Grammar Lens: two sentences that were built this way all along}]
You have been saying this pattern for a long time without being given it:

\noindent
\begin{tabularx}{\linewidth}{@{}>{\raggedright\arraybackslash}X>{\raggedright\arraybackslash}X@{}}
\toprule
\textbf{sentence} & \textbf{what it is} \\
\midrule
\hi{एक} \hi{बज} \hi{रहा} \hi{है} & \emph{baj-} bare, \textbf{\hi{रहा}}, \textbf{\hi{है}} \\
\hi{बारिश} \hi{हो} \hi{रही} \hi{है} & \emph{ho-} bare, \textbf{\hi{रही}}, \textbf{\hi{है}} \\
\bottomrule
\end{tabularx}

Neither was a lump. The clock one is masculine and the rain one is feminine, which is the only difference between them, and it is the middle word that shows it — exactly as the grid above says it should.
\end{grammarlens}

\begin{grammarlens}[title={Grammar Lens: against the habit}]
\noindent
\begin{tabularx}{\linewidth}{@{}>{\raggedright\arraybackslash}X>{\raggedright\arraybackslash}X@{}}
\toprule
\textbf{form} & \textbf{the verb} \\
\midrule
habit & \hi{बोलता} \hi{हूँ} \\
in progress & \hi{बोल} \hi{रहा} \hi{हूँ} \\
earlier habit & \hi{बोलता} \hi{था} \\
\bottomrule
\end{tabularx}

Three tenses off one stem, and the only thing you have to look at is what sits immediately after it.
\end{grammarlens}

\subsection*{Your turn}
Say these aloud:

\begin{itemize}
\raggedright
  \item all six rows of the grid
  \item the clock sentence, and name its three pieces
  \item the rain sentence, and say why its middle word ends in \textbf{-\hi{ई}}
  \item the habit and the continuous of \emph{to speak}, one after the other
\end{itemize}

\begin{tcolorbox}[breakable,colback=teal!4,colframe=teal!35!black,title={Before you move on}]
What shape is the stem in, in every row? (Bare.) Why does the rain sentence use \textbf{\hi{रही}}? (\emph{Bārish} is feminine.) What separates the habit from the continuous? (What sits after the stem.)
\end{tcolorbox}
\section[āp kyā kar rahe haiṁ?]{\hi{आप} \hi{क्या} \hi{कर} \hi{रहे} \hi{हैं}? — the question, and four ways to answer it}

\label{lesson:HI-C89-sintesis-jaari}



\begin{quote}
You have \textbf{\hi{क्या}} for \emph{what}. Put it in front of the plural continuous and you have the question this lesson is about.
\end{quote}

\subsection*{The exchange}
\begin{quote}
— \hi{आप} \hi{क्या} \hi{कर} \hi{रहे} \hi{हैं}? — \hi{मैं} \hi{हिंदी} \hi{बोल} \hi{रहा} \hi{हूँ।}
\end{quote}

Three things are at work. \textbf{\hi{क्या}} asks. \textbf{\hi{रहे} \hi{हैं}} is the plural because the question is to somebody addressed as \textbf{\hi{आप}}. The answer drops back to \textbf{\hi{रहा} \hi{हूँ}}, because the answer is about one man speaking for himself — a woman answers \textbf{\hi{बोल} \hi{रही} \hi{हूँ}}, and the question does not change, because the asker does not know yet who will answer.

\subsection*{How to answer}
Answer it four ways, out loud:

\begin{itemize}
\raggedright
  \item you are speaking Hindi
  \item you are going home
  \item you are eating
  \item and then say what you do \textbf{in general}, rather than now
\end{itemize}

The fourth one is the trap and the test. \emph{What are you doing?} takes the continuous; \emph{what do you do?} takes the habit, and the two questions sound nearly the same in English.

\begin{grammarlens}[title={Grammar Lens: what a listening paper does with this}]
\noindent
\begin{tabularx}{\linewidth}{@{}>{\raggedright\arraybackslash}X>{\raggedright\arraybackslash}X@{}}
\toprule
\textbf{what you hear} & \textbf{what it tells you} \\
\midrule
...\hi{रहा} \hi{है} & one man, right now \\
...\hi{रही} \hi{है} & one woman, right now \\
...\hi{रहे} \hi{हैं} & more than one, or \textbf{\hi{आप}} \\
\bottomrule
\end{tabularx}

The middle word is doing almost all the work, and it sits at the end of the sentence where a listener hears it last. Train the ear on those three shapes and a personal-account listening item becomes three decisions instead of one guess.
\end{grammarlens}

\subsection*{Your turn}
Say these aloud:

\begin{itemize}
\raggedright
  \item \textquotedblleft{}\hi{आप} \hi{क्या} \hi{कर} \hi{रहे} \hi{हैं}?\textquotedblright{}
  \item \textquotedblleft{}\hi{मैं} \hi{हिंदी} \hi{बोल} \hi{रहा} \hi{हूँ}\textquotedblright{}
  \item the answer a woman would give
  \item the same sentence as a habit rather than as a moment
\end{itemize}

\begin{tcolorbox}[breakable,colback=teal!4,colframe=teal!35!black,title={Before you move on}]
Why is the question \textbf{\hi{रहे} \hi{हैं}} and the answer \textbf{\hi{रहा} \hi{हूँ}}? (The question is to \textbf{\hi{आप}}; the answer is about one person.) Which word carries the gender? (The middle one.) Which form answers \emph{what do you do?} (The habit.)
\end{tcolorbox}
